\begin{helper}[Exact tree-sequence reordering]
For every Scheide structural configuration and every
\noderef{ScheideStructuralSelectionData}, there is a reordered
tree-sequence $T'$ carrying a
\noderef{ScheideStructuralReorderingData}.
\end{helper}
\begin{proof}
Write $a=D.a$, $r=D.r$, $\upsilon=D.\upsilon$, and $\rho=D.\rho$.
Define the alternating list $T'$ exactly as in
\noderef{ScheideStructuralReorderingData}: first the old base, then the
$a$ complementary blocks in $\upsilon$-order, and finally the $r$ retained
blocks in $\rho$-order.  The retained/complement partition and
$C.s=a+r$ from \noderef{ScheideStructuralSelectionData} show that these
three intervals are disjoint and exhaustive.  Together with the exact
length field of \noderef{ScheideStructuralConfiguration}, they show that
$T'$ has $T.\operatorname{length}+1$ vertices and
$T.\operatorname{length}$ edges and establish every displayed position
identity in the reordering datum.

We verify the fields of \noderef{TreeSequence}.  Positivity follows from
$C.p\ge2$.  The first exact edge identity is the unchanged old
$C.\operatorname{baseEdge}(0)=e_0$.  The vertex-position map is a
permutation of the old vertex-position map: every old base entry is in the
first interval and every old block index is in exactly one of the two
later intervals.  The identical argument for edge positions is an
identity-preserving permutation of the old edge-position map.  Hence the
old vertex and edge injectivity fields prove the new ones; in particular,
distinct parallel edge identities remain distinct.

For attachment, an old-base edge retains its old earlier endpoint.  An
absorbed edge has the form
$C.\operatorname{blockEdge}(\upsilon(h),j)$, and the block-attachment
field of \noderef{ScheideStructuralConfiguration} identifies its ends
with its displayed new vertex and some $C.y(q)$ in the initial base.
That base vertex occurs strictly earlier.  The same argument with
$\rho(t)$ proves attachment for each retained edge.  This completes the
tree-sequence construction.

The new sequence satisfies \noderef{TashkinovTree}.  The globally maximum
field of \noderef{ScheideStructuralConfiguration}, unfolded through
\noderef{GlobalMaximumTriple}, supplies this property for the old tree.
Every nonfirst old-base edge therefore retains an earlier missing-color
witness in the unchanged base prefix.  For an absorbed or retained block
edge, the block-color and missing-on-base fields give respectively its
color and a missing witness among the first $C.p$ vertices, strictly
before the block edge.  Thus every nonfirst new edge has the required
earlier witness.

The same exhaustive permutation of vertex positions and
\noderef{TreeSequenceVertexSet} give $V(T')=V(T)$.  It remains to compare
exits.  Unfold \noderef{ScheideExitVertices} and
\noderef{ScheideExitVertex}.  Given witnesses
$u,\gamma,\delta,z,P$ and a selected position for one tree, use the same
vertices, colors, exact incidence walk, and position for the other.
The equality $V(T')=V(T)$ transports start membership, the defective-color
membership after unfolding \noderef{DefectiveColors},
\noderef{DefectiveColor}, \noderef{BoundaryColorEdges}, and
\noderef{BoundaryEdges}, membership of the selected exit, and every
later-vertex exclusion.  In \noderef{ScheideDefectiveChainPath} it transports the start
and terminal inside/outside conditions; the path and full-component
conditions do not mention the ordered tree beyond its vertex set.  This
proves each implication and hence equality of the exact exit sets.  All
fields of \noderef{ScheideStructuralReorderingData} now hold.
\end{proof}
